
%%%%%%%%%%%%%%%%%%%%%%%%%%%%%%%%%%%%%%%%%%%%%%%%%%%%%%%%%%%%%%%%%%
% Sample template for MIT Junior Lab Student Written Summaries
% Available from http://web.mit.edu/8.13/www/Samplepaper/sample-paper.tex
% Last Updated April 12, 2007
% Adapted from the American Physical Societies REVTeK-4 Pages
% at http://publish.aps.org
%\newcommand{\mjmlegacy}{}

\setlength{\paperheight}{11in}
% http://tex.stackexchange.com/questions/74636/mla-package-and-thumbpdf
\makeatletter
\@namedef{ver@thumbpdf.sty}{}
\makeatother
%\documentclass[aps,secnumarabic,balancelastpage,amsmath,amssymb,nofootinbib]{revtex4}



\ifdefined\mjmlegacy
\documentclass[aps,secnumarabic,balancelastpage,amsmath,amssymb,nofootinbib]{revtex4}
\else
\documentclass[showkeys,aps,secnumarabic,balancelastpage,amsmath,amssymb,nofootinbib]{revtex4-2}
\usepackage{array}[=2016-10-06]
\fi


%http://tex.stackexchange.com/questions/119905/insert-multiple-figures-in-latex

\usepackage[nomessages]{fp}    %mjm   needed for chemfig and mol2chemfig computed angles  
\usepackage{siunitx}        %mjm  appendix table subsections  
%\usepackage{morefloats}        %mjm   saving up figs for the end   
% fking incompatibvle fuxk ing floatrow 
%\usepackage{float}        %mjm  appendix table subsections  
\usepackage{pbox}        %mjm box off junk  
\usepackage{comment}        %mjm  see the build options  
\usepackage{framed}        %mjm box off junk  
\usepackage{lgrind}        % convert program listings to a form includable in a LaTeX document
% comment out for biblatex test
% causes problems on texlive 2023+
%\usepackage{chapterbib}    % allows a bibliography for each chapter (each labguide has it's own)
%\usepackage{biblatex}   
\usepackage{color}         % produces boxes or entire pages with colored backgrounds
\usepackage{graphics}      % standard graphics specifications
\usepackage[pdftex]{graphicx}      % alternative graphics specifications
%\usepackage{graphicx}      % alternative graphics specifications
\usepackage{longtable}     % helps with long table options
\usepackage{epsf}          % old package handles encapsulated post script issues
\usepackage{bm}            % special 'bold-math' package
\usepackage{url}            % path for stupid jobname   
%\usepackage{asymptote}     % For typesetting of mathematical illustrations
\usepackage{thumbpdf}
\usepackage[colorlinks=true]{hyperref}  % this package should be added after all others
%\usepackage[draft=false, x-bib-pages=\input{\mjmbasename.last_page}, colorlinks=true]{hyperref}  % this package should be added after all others
%\usepackage[draft=false, x-bib-pages=\input{allbib.last_page}, colorlinks=true]{hyperref}  % this package should be added after all others
                                        % use as follows: \url{http://web.mit.edu/8.13}

%http://tex.stackexchange.com/questions/12676/add-notes-under-the-table
%\usepackage{booktabs,caption,fixltx2e}
% not work with subfig????
\usepackage[CaptionAfterwards]{fltpage}
\usepackage{lipsum}

% this does not  work .... 
%\usepackage[utf8]{inputenc}
%\usepackage{tabulary}
\usepackage[para,online,flushleft]{threeparttable}
%\usepackage{threeparttable}

%\usepackage{chemfig}
\usepackage[version=3]{mhchem}        %  
\usepackage{mol2chemfig}
% chemfig vriables maybe
\usepackage{xstring}        %mjm  appendix table subsections  
%\usepackage{chemformula}
% \usepackage{chemmacros}
\usepackage{floatrow}
\usepackage{fancyhdr}
%  underscore in jobmane f 
% https://latex.org/forum/viewtopic.php?t=2975
% this also messes up pdftotext  
%\usepackage[T1]{fontenc}
\usepackage{dcolumn} % https://tex.stackexchange.com/questions/2746/aligning-numbers-by-decimal-points-in-table-columns

\usepackage{catchfile} % mjmaddbib needs to read page count file
\pagestyle{fancy}




\newcolumntype{.}[1]{D{.}{.}{#1}}
%\usepackage[maxfloats=30]{morefloats}   %mjm   saving up figs for the end   
% no param on old version stuck at 36
\usepackage{morefloats}        %mjm   saving up figs for the end   
\usepackage{graphicx}        %mjm   saving up figs for the end   

% will need modificaitons 

% 2020-10-18 extract some new boilerplate 

% https://tex.stackexchange.com/questions/121601/automatically-wrap-the-text-in-verbatim
\usepackage{listings}
\lstset{
basicstyle=\small\ttfamily,
columns=flexible,
breaklines=true
}
%%%%%%%%%%%%%%%%%%%%%%%%%%%%% utilitites

\newcommand{\mjmblackbox}[2]{
 \fbox{
% thi does not ing work right 
\begin{minipage}[t]{\textwidth}
{ \centering{\bf{#1 : }} }
\par
#2
\end{minipage}
}
}
\newcommand{\mjmblackboxno}[2]{
 \fbox{
% thi does not ing work right 
\begin{minipage}[t]{\textwidth}
{ \centering{\bf{#1 : }} }
#2
\end{minipage}
}
}





%%%%%%%%%%%%%%%%%%%%%%%% biblio stuff



\newcommand{\checkrel}[1]{%
  \ifcsname#1\endcsname%
\newcommand{\mjmstatus}{  public NOTES }
\newcommand{\mjmversion}{\mjmrelease}
  \else%
\newcommand{\mjmstatus}{ NOT public NOTES }
\newcommand{\mjmversion}{0.00}
  \fi%
}




% https://tex.stackexchange.com/questions/18089/are-there-any-command-for-producing-the-bibtex-logo
%\def\BibTeX{{\rm B\kern-.05em{\sc i\kern-.025em b}\kern-.08em
%    T\kern-.1667em\lower.7ex\hbox{E}\kern-.125emX}}

\newcommand{\biblogo}{
{{\rm B\kern-.05em{\sc i\kern-.025em b}\kern-.08em
    T\kern-.1667em\lower.7ex\hbox{E}\kern-.125emX}}
{ }  }
%\newcommand{\biblogo}{ Bibte{\it X}  { }  }
\newcommand{\latexlogo}{ \LaTeX  { }  }
\newcommand{\bomtexlogo}{ BomTe{\it X}  { }  }






\newcommand{\mjmvirus}{SARS-Cov-2 }
\newcommand{\mjmdisease}{covid-19 }
\newcommand{\Mjmdisease}{Covid-19 }
\newcommand{\mjmlogo}{ MUQED { }  }
\newcommand{\mjmlinkedin}{ {\bf LinkedIn} { }  }








% https://tex.stackexchange.com/questions/121601/automatically-wrap-the-text-in-verbatim
\usepackage{listings}
\lstset{
basicstyle=\small\ttfamily,
columns=flexible,
breaklines=true
}

% none of this fking fking works for a f F 
\newcommand{\mjmverbatim}{lstlisting}
\newcommand{\mjmbeginverbatim}{\begin{lstlisting}}
\newcommand{\mjmendverbatim}{\end{lstlisting}}

\newcommand{\mjmmangle}[1]{keep/#1}




\usepackage[section]{placeins}


%http://tex.stackexchange.com/questions/119905/insert-multiple-figures-in-latex

\usepackage[nomessages]{fp}    %mjm   needed for chemfig and mol2chemfig computed angles  
\usepackage{siunitx}        %mjm  appendix table subsections  
%\usepackage{morefloats}        %mjm   saving up figs for the end   
% fking incompatibvle fuxk ing floatrow 
%\usepackage{float}        %mjm  appendix table subsections  
\usepackage{pbox}        %mjm box off junk  
\usepackage{comment}        %mjm  see the build options  
\usepackage{framed}        %mjm box off junk  
\usepackage{lgrind}        % convert program listings to a form includable in a LaTeX document
% comment out for biblatex test
% causes problems on texlive 2023+
%\usepackage{chapterbib}    % allows a bibliography for each chapter (each labguide has it's own)
%\usepackage{biblatex}   
\usepackage{color}         % produces boxes or entire pages with colored backgrounds
\usepackage{graphics}      % standard graphics specifications
\usepackage[pdftex]{graphicx}      % alternative graphics specifications
%\usepackage{graphicx}      % alternative graphics specifications
\usepackage{longtable}     % helps with long table options
\usepackage{epsf}          % old package handles encapsulated post script issues
\usepackage{bm}            % special 'bold-math' package
\usepackage{url}            % path for stupid jobname   
%\usepackage{asymptote}     % For typesetting of mathematical illustrations
\usepackage{thumbpdf}
\usepackage[colorlinks=true]{hyperref}  % this package should be added after all others
%\usepackage[draft=false, x-bib-pages=\input{\mjmbasename.last_page}, colorlinks=true]{hyperref}  % this package should be added after all others
%\usepackage[draft=false, x-bib-pages=\input{allbib.last_page}, colorlinks=true]{hyperref}  % this package should be added after all others
                                        % use as follows: \url{http://web.mit.edu/8.13}

%http://tex.stackexchange.com/questions/12676/add-notes-under-the-table
%\usepackage{booktabs,caption,fixltx2e}
% not work with subfig????
\usepackage[CaptionAfterwards]{fltpage}
\usepackage{lipsum}

% this does not  work .... 
%\usepackage[utf8]{inputenc}
%\usepackage{tabulary}
\usepackage[para,online,flushleft]{threeparttable}
%\usepackage{threeparttable}

%\usepackage{chemfig}
\usepackage[version=3]{mhchem}        %  
\usepackage{mol2chemfig}
% chemfig vriables maybe
\usepackage{xstring}        %mjm  appendix table subsections  
%\usepackage{chemformula}
% \usepackage{chemmacros}
\usepackage{floatrow}
\usepackage{fancyhdr}
%  underscore in jobmane f 
% https://latex.org/forum/viewtopic.php?t=2975
% this also messes up pdftotext  
%\usepackage[T1]{fontenc}
\usepackage{dcolumn} % https://tex.stackexchange.com/questions/2746/aligning-numbers-by-decimal-points-in-table-columns

\usepackage{catchfile} % mjmaddbib needs to read page count file
\pagestyle{fancy}




\newcolumntype{.}[1]{D{.}{.}{#1}}
%\usepackage[maxfloats=30]{morefloats}   %mjm   saving up figs for the end   
% no param on old version stuck at 36
\usepackage{morefloats}        %mjm   saving up figs for the end   
\usepackage{graphicx}        %mjm   saving up figs for the end   

% will need modificaitons 
%mjmlegacy
%
% 2020-10-18 extract some new boilerplate 

% https://tex.stackexchange.com/questions/121601/automatically-wrap-the-text-in-verbatim
\usepackage{listings}
\lstset{
basicstyle=\small\ttfamily,
columns=flexible,
breaklines=true
}
%%%%%%%%%%%%%%%%%%%%%%%%%%%%% utilitites

\newcommand{\mjmblackbox}[2]{
 \fbox{
% thi does not ing work right 
\begin{minipage}[t]{\textwidth}
{ \centering{\bf{#1 : }} }
\par
#2
\end{minipage}
}
}
\newcommand{\mjmblackboxno}[2]{
 \fbox{
% thi does not ing work right 
\begin{minipage}[t]{\textwidth}
{ \centering{\bf{#1 : }} }
#2
\end{minipage}
}
}





%%%%%%%%%%%%%%%%%%%%%%%% biblio stuff



\newcommand{\checkrel}[1]{%
  \ifcsname#1\endcsname%
\newcommand{\mjmstatus}{  public NOTES }
\newcommand{\mjmversion}{\mjmrelease}
  \else%
\newcommand{\mjmstatus}{ NOT public NOTES }
\newcommand{\mjmversion}{0.00}
  \fi%
}




% https://tex.stackexchange.com/questions/18089/are-there-any-command-for-producing-the-bibtex-logo
%\def\BibTeX{{\rm B\kern-.05em{\sc i\kern-.025em b}\kern-.08em
%    T\kern-.1667em\lower.7ex\hbox{E}\kern-.125emX}}

\newcommand{\biblogo}{
{{\rm B\kern-.05em{\sc i\kern-.025em b}\kern-.08em
    T\kern-.1667em\lower.7ex\hbox{E}\kern-.125emX}}
{ }  }
%\newcommand{\biblogo}{ Bibte{\it X}  { }  }
\newcommand{\latexlogo}{ \LaTeX  { }  }
\newcommand{\bomtexlogo}{ BomTe{\it X}  { }  }






\newcommand{\mjmvirus}{SARS-Cov-2 }
\newcommand{\mjmdisease}{covid-19 }
\newcommand{\Mjmdisease}{Covid-19 }
\newcommand{\mjmlogo}{ MUQED { }  }
\newcommand{\mjmlinkedin}{ {\bf LinkedIn} { }  }








% https://tex.stackexchange.com/questions/121601/automatically-wrap-the-text-in-verbatim
\usepackage{listings}
\lstset{
basicstyle=\small\ttfamily,
columns=flexible,
breaklines=true
}

% none of this fking fking works for a f F 
%mjmlegacy
%\newcommand{\mjmverbatim}{lstlisting} \newcommand{\mjmbeginverbatim}{\begin{lstlisting}} \newcommand{\mjmendverbatim}{\end{lstlisting}} \newcommand{\mjmmangle}[1]{keep/#1}



%# CHANGE VERSION AND STATUS MANUALLY 
% need a draft/notes/release flag

% https://tex.stackexchange.com/questions/5894/latex-conditional-expression
%At the command-line, you can do \def\MYFLAG{} and then test if \MYFLAG is defined in your document (or an included style file) with \ifdefined\MYFLAG ... \else ... \fi.
% needs trailing space for the sample bibtex doh
% leading spaces mess up the entry thought 
\def\xxmjmrelease{0.10 }
\ifdefined\mjmrelease
\newcommand{\mjmstatus}{ PUBLIC NOTES }
\newcommand{\mjmversion}{\mjmrelease} %%%%%%%%%%%%%
\newcommand{\mjmtrno}{MJM-2026-002}
\newcommand{\mjmbib}{\mjmtrno-\mjmversion}
\newcommand{\mjmstatuswarn}{{\bf{  }}   }
% 2021-09-29 wanted version wth for brownie
%\newcommand{\mjmbib}{\mjmtrno-\mjmversion-\mjmrelease}
%\newcommand{\mjmbib}{\mjmtrno}
\else
\newcommand{\mjmstatus}{ NOT public NOTES }
\newcommand{\mjmversion}{0.00} %%%%%%%%%%%%%
\newcommand{\mjmtrno}{MJM-2026-002}
%\newcommand{\mjmbib}{\mjmtrno-\mjmversion}
\newcommand{\mjmbib}{\mjmtrno}
%\newcommand{\mjmstatuswarn}{  }
\newcommand{\mjmstatuswarn}{{\bf{This document is a non-public DRAFT and contents may be speculative or undocumented or simple musings and should be read as such.  }}   }
\fi

%\newcommand{\mjmstatus}{ NOT public NOTES }



\newcommand{\mjmtitle}{Sildenafil Seems Beneficial in Dog with Heartworm 8 years ago }
\newcommand{\mjmsubtitle}{Ask Your Veterinanarian about Sildenafil to keep heart from beating like a big bass drum : The Happy Pill}
\begin{comment}
\newcommand{\xxmjmmakedate}{ }
\newcommand{\xxmjmauthor}{Mike J Marchywka }
\newcommand{\xxmjmbasename}{\jobname}
\newcommand{\xxmjmaddbio}{mjm_tr,releases}
\newcommand{\xxmjmversion}{ 0.00 }
\newcommand{\xxmjmtrno}{MJM-2018-009}
\newcommand{\xxmjmbibday}{12}
\newcommand{\xxmjmbibmo}{12}
\newcommand{\xxmjmbibyear}{2018}
\newcommand{\xxmjmemail}{marchywka@hotmail.com}
\end{comment}



\ifdefined\mjmlegacy
\newcommand{\expandableinput}[1]{\input{#1}}
\else
% https://tex.stackexchange.com/questions/583927/misplaced-noalign-error-with-input-in-a-table-after-the-2020-fall-latex-releas/583939#583939
\ExplSyntaxOn % providing \expandableinput
\cs_new:Npn \expandableinput #1
  { \use:c { @@input } { \file_full_name:n {#1} } }
\ExplSyntaxOff
\fi



\ifdefined\grad
\else
\newcommand{\grad}{\nabla}
\fi
\newcommand{\mjmquote}[1]{{\centering {\it #1 }}}
\newcommand{\mjmwisdom}[1]{{\centering {\it #1 }}}
\newcommand{\laplace}{\nabla^{2}}
\newcommand{\mjmdx}[2]{\left(\frac{\partial #1 }{\partial #2} \right) }
\newcommand{\mjmdxop}[2]{\frac{\partial  }{\partial #2}\left( #1 \right) }
\newcommand{\mjmdxdx}[2]{\left(\frac{\partial^{2} #1 }{\partial #2^{2}} \right) }
\newcommand{\mjmdxo}[2]{\frac{\partial #1 }{\partial #2} }
\newcommand{\mjmdxx}[2]{\left(\frac{\partial^{2} #1 }{\partial {#2}^{2}} \right) }
\newcommand{\mjmD}[3]{\frac{d^{#1} #2 }{d{#3}^{#1}}  }
\newcommand{\mjmDnt}[2]{\frac{d^{#1} #2 }{dt^{#1}}  }
\newcommand{\mjmDparen}[2]{ {#1}^{(#2)} }
\newcommand{\mjmdxy}[3]{\left(\frac{\partial^{2} #1 }{\partial {#2}\partial{#3}} \right) }
\newcommand{\mjmdxyn}[5]{\left(\frac{\partial^{#1} }{\partial^{#2} {#3}\partial^{#4}{#5}} \right) }
\newcommand{\mjmdsq}[2]{\left(\frac{\partial #1 }{\partial #2} \right) ^{2}}
\newcommand{\mjmupdated}[2]{\p  Updated on $#1$ from source #2 \p }
% see tug mail archives, based on discussion 
%\bool_new:N \l_tmpa_bool
\newif\ifbibstarted
\newif\ifbibnamed
\bibstartedfalse
\bibnamedfalse
\newcommand{\mjmtotalbib}{}
%\def\foo{}
\newcommand{\mjmsummabib}[2]{
%\renewcommand{\mjmtotalbib}{ \mjmtotalbib, #1 = #2 }
%\let\foo{ #1 = #2}
}
%%%%%%%%%%%%%%%%%%%%%%%%%%%%%%%%%%%%%%%%%%%%%%%%%%%%%%%%%%%%%%%%

\iffalse
@software{,
  author = {Michael J Marchywka},
  city = {Jasper GA 30143 USA},
  title = { A one-file library for adding machine and human readable bibtex to an article },
abstract={ A simple include file to make bibtex available 
in a document in both machine and human readable format. 
Machine readable is added to extended information 
which can be read with tools such as exiftool ( https://exiftool.org/ ).
While typically not complete at time of pdf creation,
other tools such as toobib can be used to complete the citation
and of course publishers may be able to modify it too as 
more is known. The human readable form need not beincluded
in document types not suited for that but then automated citation
may still be easy.  
See also some exchanges on the Texhax mailing list @tug.org },
institution={},
license={Knowledge sir should be free to all },
publisher={Mike Marchywka},
email={marchywka@hotmail.com},
authorid={orcid.org/0000-0001-9237-455X},
  filename = {mjmaddbib.tex},
  url = {},
  version = {0.0.0},
  date-started = {}
}

<one line to give the program's name and a brief idea of what it does.>


Conceived and written by Mike Marchywka from 2019 to present.
See dates in individual code pieces as they were 
generated from my wizards. 
Copyright (C) <year> <name of author>


This program is free software: you can redistribute it and/or modify it under
the terms of the GNU General Public License as published by the Free Software
Foundation, either version 3 of the License, or (at your option) any later
version.

This program is distributed in the hope that it will be useful, but WITHOUT ANY
WARRANTY; without even the implied warranty of  MERCHANTABILITY or FITNESS FOR
A PARTICULAR PURPOSE. See the GNU General Public License for more details.

You should have received a copy of the GNU General Public License along with
this program.  If not, see <http://www.gnu.org/licenses/>.

   THIS SOFTWARE IS PROVIDED BY THE COPYRIGHT HOLDERS AND CONTRIBUTORS
   "AS IS" AND ANY EXPRESS OR IMPLIED WARRANTIES, INCLUDING, BUT NOT
   LIMITED TO, THE IMPLIED WARRANTIES OF MERCHANTABILITY AND FITNESS FOR
   A PARTICULAR PURPOSE ARE DISCLAIMED.  IN NO EVENT SHALL THE COPYRIGHT OWNER OR
   CONTRIBUTORS BE LIABLE FOR ANY DIRECT, INDIRECT, INCIDENTAL, SPECIAL,
   EXEMPLARY, OR CONSEQUENTIAL DAMAGES (INCLUDING, BUT NOT LIMITED TO,
   PROCUREMENT OF SUBSTITUTE GOODS OR SERVICES; LOSS OF USE, DATA, OR
   PROFITS; OR BUSINESS INTERRUPTION) HOWEVER CAUSED AND ON ANY THEORY OF
   LIABILITY, WHETHER IN CONTRACT, STRICT LIABILITY, OR TORT (INCLUDING
   NEGLIGENCE OR OTHERWISE) ARISING IN ANY WAY OUT OF THE USE OF THIS
   SOFTWARE, EVEN IF ADVISED OF THE POSSIBILITY OF SUCH DAMAGE.


\fi

%\usepackage[pdftex]{graphicx}      % alternative graphics specifications
\usepackage{hyperref}      %


\newcommand{\mjmstartbib}[2]
{
\def\mjmbibentry{@#1\{#2}
\def\mjmbiboneentry{@#1\{#2}
\def\mjmbibpre{x-bib{-}}
\def\mjmday{\day}
%\mjmaddbib{run-day}{\mjmday}
%\mjmaddbib{run-month}{\expandafter\month}
%\mjmaddbib{run-year}{\year}
\mjmaddbib{filename}{\jobname}
\mjmaddbib{run-date}{\today}
}
% modify toobib and change this from "x-bib" to something like bibtex or mjm- whatever 
\newcommand{\mjmbibmunge}[1] {x-bib-#1}

\newcommand{\mjmaddbib}[2]
{
{\hypersetup{pdfinfo={{\mjmbibmunge{#1}}={#2}}}}
\edef\mjmbibentry{\unexpanded\expandafter{\mjmbibentry,}

    #1 =\{#2\} }
%\edef\mjmbiboneentry{\unexpanded\expandafter{\mjmbiboneentry,}
% this works but gives warning and linebreak is removed.
%\edef\mjmbiboneentry{\unexpanded\expandafter{\mjmbiboneentry,\linebreak}
\edef\mjmbiboneentry{\unexpanded\expandafter{\mjmbiboneentry, }
#1 =\{#2\}
}

} % mjmaddbib


\newcommand{\mjmshowbib}
{

\mjmbibentry

\}
} % mjmshowbib

\newcommand{\mjmshowbibone}
{
\mjmbiboneentry \}
} % mjmshowbibone

\newcommand{\mjmdonebib}
{

{\hypersetup{pdfinfo={\mjmbibmunge{bibtex}={\mjmshowbibone}}}}

} % mjmdonebib


% af 
\newcommand{\mjminputlisting}[2]
{
%\begin{figure}[H]
\lstinputlisting{#1}
%\caption{#2}


#2

%\end{figure}
}

%%%%%%%%%%%%%%%%%%%%%%%%%%%%%%%%%%%%%%%%%%%%%%%%%%%%%%%%%%%%%%%%%%

%\newcommand{\mjmaddbib}[2]{\hypersetup{ pdfinfo={ x-bib-#1 = {#2}}}\mjmsummabib{#1}{#2}}

\newcommand{\mjmaddbibonly}[2]{\hypersetup{ pdfinfo={ x-bib-#1 = {#2}}}}
\newcommand{\mjmaddbibe}[2]{\hypersetup{ pdfinfo= x-bib-#1 = #2}}

\newcommand{\mjmtable}[2]{
\begin{table}[H] \centering
\begin{tabular}{#1}
#2
\end{tabular}
\end{table}

}
% for archiving and file list 
% David Carlisle You can use   \textbf{\detokenize{ pmg_ratios.svg}} \IfFileExists{ pmg_ratios.sv}{yes}{no} to make the tokens safe for typesetting in text mode.


\newcommand{\mjmusesitem}[2]{

%\item {\bf $ #1 $  } : #2  % 1 is #1 and 2 is  #2 xxx 
\item {{\detokenize{ #1 } }} : #2  % 1 is #1 and 2 is  #2 xxx 
\IfFileExists{#1}{}{{\bf not found} }

} % mjmuses

\newcommand{\mjmreleasewarning}
{
{\bf This is a draft and has not been peer reviewed or completely proof 
read but released in some state where it seems worthwhile given 
time or other constraints. Typographical errors are quite likely 
particularly in manually entered numbers. This work may 
include  output from software which has not been fully debugged.  
For information only, not for use for any particular purpose see
fuller disclaimers in the text.  Caveat Emptor.}
}

\newcommand{\mjmwarningtoo}
{\bf {This is a draft which may not have been fully proofread and certainly
not peer reviewed. Read the disclaimers and take them seriously.
The reader is assumed familiar with the related literature and
controversial issues. For information and thought only not intended
for any particular purpose. Caveat Emptor  }}

\newcommand{\mjmwarntopic}
{\bf {  This work addresses a controversial topic and likely advances one
or more viewspoints that are not well accepted in an attempt to
resolve confusion.   
The reader is assumed familiar with the related literature and
controversial issues and in any case should seek additional 
input from sources the reader trusts likely with differing opinions. 
For information and thought only not intended
for any particular purpose. Caveat Emptor  }}


\newcommand{\mjmwarnfeed}
{\bf { Note that any item given to a non-human must be checked for safety alone and in combination with other ingredients or medicines  for that animal. Animals including dogs and cats have decreased tolerance for many common ingredients in things meant for human consumption. }}


\newcommand{\mjmwarnme}
{\bf { I am not a veterinarian or a doctor or health care professional
 and this is not particular advice
for any given situation.  Read the disclaimers in the appendicies or text, take them seriously and take prudent steps 
to evaluate this information. 
 }}

\newcommand{\mjmexplainbib}
{{The release may use an experimental bibliography code that
is not designed to achieve a particular format but to
allow multiple links to reference works with modifications
to the query string to allow identification of the citing
work for tracking purposes. This may be useful for a bill-of-materials
and purchases later.
}}








\newcommand{\mjmauthor}{Mike J Marchywka }
\newcommand{\mjmmakedate}{2026-07-09 }
\newcommand{\mjmbasename}{\jobname}
\newcommand{\mjmaddbio}{mjm_tr,releases}
%\newcommand{\mjmversion}{0.00}
%\newcommand{\mjmtrno}{MJM-2026-002}
%\newcommand{\mjmbibday}{09}
%\newcommand{\mjmbibmo}{07}
%\newcommand{\mjmbibyear}{2026}

% the build script changes these to creation day doh 
\newcommand{\mjmbibday}{09}
\newcommand{\mjmbibmo}{07}
\newcommand{\mjmbibyear}{2026}


\newcommand{\mjmmakebibday}{\number\day}
\newcommand{\mjmmakebibmo}{\number\month}
\newcommand{\mjmmakebibyear}{\number\year}

\newcommand{\mjmbibtype}{techreport}

\newcommand{\mjmbibname}{marchywka-\mjmbib}
\mjmstartbib{\mjmbibtype}{\mjmbibname}


\newcommand{\mjmemail}{marchywka@hotmail.com}
%\newcommand{\mjmaddr}{157 Zachary Dr Talking Rock GA 30175 USA}
\newcommand{\mjmaddr}{157 Zachary Dr Talking Rock GA 30175 USA}
\mjmaddbib{title}{\mjmtitle}
\mjmaddbib{subtitle}{\mjmsubtitle}
\mjmaddbib{author}{\mjmauthor}
\mjmaddbib{type}{\mjmbibtype}
%\mjmaddbib{name}{marchywka-\mjmbib}
\mjmaddbib{name}{\mjmbibname}
\mjmaddbib{number}{\mjmtrno}
\mjmaddbib{version}{\mjmversion}
\mjmaddbib{institution}{not institutionalized, independent }
\mjmaddbib{address}{ \mjmaddr}
\mjmaddbib{date}{\today}
\mjmaddbib{startdate}{\mjmbibyear -\mjmbibmo -\mjmbibday }
%\mjmaddbib{day}{\mjmbibday}
%\mjmaddbib{month}{\mjmbibmo}
%\mjmaddbib{year}{\mjmbibyear}
\mjmaddbib{day}{\mjmmakebibday}
\mjmaddbib{month}{\mjmmakebibmo}
\mjmaddbib{year}{\mjmmakebibyear}

\mjmaddbib{author1email}{\mjmemail}
\mjmaddbib{contact}{\mjmemail}
\mjmaddbib{author1id}{orcid.org/0000-0001-9237-455X}
\CatchFileEdef\mjmpages{\mjmbasename.last_page}{\endlinechar=-1\relax}
% TODO FIXME add this to the skeleton text 
%\mjmaddbib{pages}{ \input{\mjmbasename.last_page}}
\mjmaddbib{pages}{ \mjmpages}
%\mjmaddbib{filename}{\mjmbasename}

\begin{comment}
\newcommand{\checkfortwo}[1]{%
  \ifcsname#1\endcsname%
  VERSION =        \{\mjmversion \today \mjmstatus \},
  \else%
  VERSION =        \{\mjmversion\},
  \fi%
}

\end{comment}
%\mjmaddbib{bibtex}{\mjmfullbib}

\mjmdonebib



\lhead{\mjmauthor,  \mjmtrno }


%\lhead{M Marchywka,  \mjmtrno }
%\rhead{ \mjmversion not for public release}
%\rhead{ { \today }  v. \mjmversion for release without review }
%\rhead{ { \today }  v. \mjmversion NOT public DRAFT }
%\rhead{ { \today }  v. \mjmversion { }  NOT public NOTES }
\rhead{ { \today }  v. \mjmversion { }  \mjmstatus }

\newfloatcommand{capbtabbox}{table}[][\FBwidth]


%%%% build flags 
%\newlength{\desttabw}  \setlength{\desttabw}{4in}
\newlength{\desttabw}  \setlength{\desttabw}{\textwidth}
\newlength{\chainwidth}  \setlength{\chainwidth}{.4\textwidth }
\newlength{\slantwidth}  \setlength{\slantwidth}{.2\textwidth }
\newlength{\subfigwidth}  \setlength{\subfigwidth}{.3\textwidth }
\newlength{\fullfigwidth}  \setlength{\fullfigwidth}{.8\textwidth }
\newlength{\subwfigwidth}  \setlength{\subwfigwidth}{.75\textwidth }
\newlength{\subwfigwidthrot}  \setlength{\subwfigwidthrot}{\textwidth }
\newlength{\myboxwidth}  \setlength{\myboxwidth}{.3\textwidth }
\newlength{\picwidth}  \setlength{\picwidth}{.4\textwidth }
% set to center for nowmal output 
\newcommand{\destflushtab}{flushleft}
\includecomment{mdpicomment}
\excludecomment{draftcomment}
\excludecomment{badmathcomment}
\excludecomment{showworkcomment}
% this does not fing work ... 

\newcommand{\mjmed}[1]{
%\begin{mjmedx} 
[ mjm : #1   ]
%\end{mjmedx} 
}  

% thinking outload
\newcommand{\mjmtolx}[1]{}
\newcommand{\mjmtolxx}[1]{}
\newcommand{\mjmtol}[1]{
 \fbox{  
% thi does not ing work right 
\begin{minipage}[t]{\textwidth}
{ \centering{\bf{Thinking outloud}} }
\par   
#1 
\end{minipage} 
}
}

\newcommand{\mjmpicture}[3]
{
\begin{figure}[H]
{ \includegraphics[height=3in,width=4in]{keep/#1} }
\caption{#2}
\label{fig:#3}
\end{figure}
} % mjmpicture


\newcommand{\mjmaside}[1]{
 \fbox{  
% thi does not ing work right 
\begin{minipage}[t]{\textwidth}
{ \centering{\bf{Aside: }} }
\par   
#1 
\end{minipage} 
}
}




\newcommand{\mjmgraphics}[1]{#1 }
\newcommand{\mjmfullplot}[1]{\includegraphics[width=\fullfigwidth]{#1}}
%\newcommand{\mjmincludeplot}[1]{\includegraphics[width=\fullfigwidth]{#1}}
%\newcommand{\mjmincludeplot}[1]{\includegraphics[width=\subfigwidth]{#1}}
\newcommand{\mjmincludeplot}[1]{\includegraphics[height=3in,width=\fullfigwidth]{#1}}

% include here as likely to be doc specific 
%\newcommand{\mjmreffig}[1]{Fig. \ref{#1}}
\newcommand{\mjmreffig}[1]{Fig. \ref{fig:#1}}
\newcommand{\mjmreftab}[1]{Table  \ref{tab:#1}}
\newcommand{\mjmrefeqn}[1]{Eqn  \ref{eqn:#1}}


%cp yyy2.pdf ~/d/latex/keep/pp20171124biotin.pdf
\newcommand{\mjmdatedplot}[1] 
{ \includegraphics[height=3in,width=\fullfigwidth]{keep/pp20171124#1} }
% right now there are too many figs to save for the end apaprently
% hard limit is 36
\newcommand{\mjmbeginfigure}{\begin{figure}[H] }
%\newcommand{\mjmbeginfigure}{\begin{figure}[p] }
\newcommand{\mjmfigure}[1]{
%\begin{figure}[H]
\mjmbeginfigure

#1 

\end{figure}
}
%\extrafloats{100}

\newcommand{\mjmlisting}[1]
{
\begin{lstlisting} 
#1 
\end{lstlisting}
}

\newcommand{\mjmold}[1]{ } 



\newcommand{\mjmeqn}[1]{\begin{equation} #1 \end{equation}  } 





\begin{document}

\title{\mjmtitle}

\mjmsubtitle

\author         {Mike Marchywka}
\email          {\mjmemail}
\thanks{ to cite  or credit this work, see bibtex in \ref{appendix:citing} } 
\date{\today}
\affiliation{\mjmaddr}

\mjmblackboxno{Release Notes  xxxx-xx-xx : }{
not ready at all

{\bf ToDo : } Most issues common with early large manuscript.

\mjmreleasewarning

\mjmstatuswarn

\mjmwarnfeed

\mjmwarnme

\mjmwarntopic

\mjmexplainbib

}

\begin{abstract}
This work describes the results giving sildenafil to one dog, 
Happy, who has had persistent cough since recovering from
heartworm in 2019. 
She has been the topic of several recent manuscripts as
her diet is adjusted to try to remove persistent cough.
After an acute episode  of apparent heart failure, she was treated with
antibiotics and furosemdie before switching to sildenafil on
request. Her recovery appears complete from the acute episode
and her cough reduction is the most robust seen so far with even
morning cough on initial wake up become reduced. Its likely
her right sided heart and vasculature never healed properly
from the initial heartworm infection. Some common medicines like
pimodendan with positive inotropic effects may increase the
problem as additional pressure results in more pathological
hypertrophy or other adaptations. Sildenafil may improve
pulmonary dilation and have a negative inotropic effect
which may be therapeutic compared to pimobendan and ACE inhibitors
for some dogs with heartworm related symptoms. 
Importantly however, she experienced an increase in coughing
over 2 day which led to discontinuation of sildenafil due to
concerns about worsening collapsed trachea.
It may be worth consideration when heartworm damage is suspected
without any echocardiography to confirm PAH but also to consider
possibility of drug related side effects like trachea collapse.
\end{abstract}

\maketitle
\tableofcontents
\newpage


\section{Introduction  }

Heart failure is fairly common in dogs due to many
reasons. Recently, new cases of DCM emerged in breeds that
were not thought to be predisposed to the condition.
Treatment for left sided failure such
as degenerative mitral valve disease
commonly consists of drugs such as furosemide,
pimobendan, and an ACE inhbitor
\cite{Saetang_Surachetpong_Short_term_effects_2020}
. 
Pimobendan shows some early signs of beneficial left sided  remodelling in
myxomatous mitral valve disease (MMVD) 
\cite{PMC13529528} but here the question is considered of how this
translates into heartworm related damage. 
Heartworm disease is another cause  and is likely responsible
for the chronic heart condition experienced by the dog
discussed here, Happy. Her story has been detailed in several
prior reports

but briefly she arrived here with heartworm in late 2018 and 
was treated with ivermection and doxycycline successfully.
She briefly stopped coughing in early 2019 but continued
to cough at varying amounts up until present in 2026.
In response to acute episodes in the past, she was prescribed
variations of the above with a diuretic, ACE inhbitors, and vetmeden.

In late 2020, another dog here, Andy, had partially recovered
from severe heartworm disease after receiving Diroban on 
2020-08-04 and 2020-09-08 but suffered from excessive abdominal
fluid accumulation.
He adopted out and, based on second hand information,  a specialist
diagnosed  PAH based on echo and prescibed sildenafil  which was thought to work well.  
Echo alone may not be complete  
\cite{TopyaPutowska_Tomaszewski_Wysokiski_Echocardiography_Pulmonary_Arterial_2021}
but given the heartworm history is probably quite convincing.
On further investigation of heartworm specific heart damage
it became clear that this may be a better approach for Happy
than the prior drugs. 
\subsection{Heartworm}

The damage caused by canine heartworm infection has
been described as, "
characterized by endothelial proliferation,
perivascular inflammation and thrombo-embolism, all
leading to chronic pulmonary hypertension and right-side
congestive heart disease. 
"
\cite{Carreton_Grandi_Morchon_Myocardial_damage_dogs_2012}.
Little change in pulmonary parameters were noted during
and 1 month after a 90 day heartworm treatment 
\cite{SerranoParreno_Carreton_CaroVadillo_Pulmonary_hypertension_dogs_2017}.
And indeed in the case of Happy little consistent change was observed
in 7 years. 



One well known target for this condition is PDE5 which 
appears to control nitric oxide levels in pulmonary
vasodilation
\cite{Ghofrani_PepkeZaba_Barbera_Nitric_oxide_pathway_2004}.
. Both PDE5 and PDE3 expression levels can increase in response
to hypoxia \cite{https://doi.org/10.1038/sj.bjp.0704984}.
Ventricular PDE5 contributes to remodelling  in mice 
\cite{Verbeken_Sipido_Herijgers_Ventricular_Phosphodiesterase_Expression_2009}.

Some forms of pulmonary artery relaxation can be inhibited by heartworm
infection although nitroglycerin response is intact even as some
doubt about the mechanisms exist
\cite{Mupanomunda_Williams_Mackenzie_Dirofilaria_immitis_heartworm_1997}.

\subsection{Right Ventricle Pulmonary Artery Death Spiral}

Pylmonary hypertension can occur for a variety of reasons
\cite{Akabane_Sato_Sakatani_Pharmacokinetics_single_dose_2020}
although the immediate interest is PAH due to heartworm damage.

The two components contributing to PAH are several structural changes
of muscle and matrix proteins as well as unregulated vasoconstriction 
\cite{Lyle_Davis_Brozovich_Regulation_Pulmonary_Vascular_2017}.


Administration of sildenafil to hypoxic rats was able to reduce
and reverse some pulmonary muscularization
\cite{Sebkhi_Strange_Phillips_Phosphodiesterase_Type_2003}/

The use of pimobendan and sildenafil in managing active
adult heartworm infection has been discussed at least since 2019
where movement of worm from the heart was observed
\cite{Tjostheim_Kellihan_Grint_Effect_sildenafil_2019}. 
In that work, the positive inotropy of pimobendan was considered benefical
in the acute situation of immediate concern. 
Simiarly, another work comparing pimobendan and sildenafil in
chemically created PH found some benefits from both or in combination  
\cite{Nakata_Tanaka_Yoshiyuki_Effects_Single_Drug_2015}.
In a chronic setting, without resolution of PAH, it is essentially
a right outflow obstruction. 
While beneficial in long term human heart failure studies, some 
excess deaths were noted in  pimobendan treatment groups although not reaching
statistical significance
\cite{Fuentes_pimobendan_2004}.

\subsection{Sildenafil }


It is available off-label in dogs and therefore
specific prescribing information is not available.
Literature reports suggest pharmacokinetics in dogs
is highly variable in dogs with natural PH
due to health state, other stomach
contents, and hepatic metabolism which may saturate over
the relevant oral dosage range of .5 to 3mg/kg
( typically 2-3 times per day )
\cite{Yata_DeFrancesco_Bonagura_Population_Pharmacokinetics_Sildenafil_2026}
. 
Earlier work with induced PH from artificial empoli
suggested increased half life at higher doses
such as 4mg/kg versus 1mg/kg
\cite{Akabane_Sato_Sakatani_Pharmacokinetics_single_dose_2020}.
Dosing in the fed state decreased peak plasma concentration and AUC
and again the non-proportional results were observed between 1 and 4mg/kg
\cite{PMID29352474}.
Accidental exposure of up to 14mg/kg have been survived according
to informal reports \cite{agraham_Sildenafil_Exposure_2026}.
Although daily dosing for 4months at 14.3mg/kg in PDE6A heterozygous
dogs did produce reversible retina abnormalities and reductions
in low light vision
\cite{KEPierce_Bartoe_Hauptman_Sildenafil_Reduces_Retinal_Sensitivity_2009}
.
Although interestingly, rat pups recovered better from hypoxic 
retinal damage with sildenafil  doses about 10mg/kg
\cite{Jung_Johnstone_Khoja_Sildenafil_Improves_Functional_2016}.
Concerns do exist about dangerout decreases in coronary artery
flow when sildenafil is combined with nitrates
\cite{Ishikura_Beppu_Hamada_Effects_Sildenafil_Citrate_2000}.


In dogs with active heartworm infection that has migrated to the
tricuspid valve, combined sildenafil and pimobendan 
in dogs without hypotension can cause worm to migrate to pulmonary artery
\cite{TJOSTHEIM201996}. 

In small breed dogs with various causes of PH the benefits of 
sildenafil were less clear when most dogs displayed syncope as well
as cough  
\cite{Johnson_Stern_Clinical_features_outcome_2020}.

One potentially serious side effect that may be relevant to
dogs is the observation of desaturation in 2 of 269 pediatric
patients leading to discontinuation of   sildenafil 
and may be related to tracheobronchomalacia 
\cite{Hansen_Airway_spasm_collapse_2019}.
Effects related to obstructive sleep apnea are being explored
although currently confusing
\cite{Andersen_Tufik_Effects_Sildenafil_2018}.

It also has a variety of effects on TRPC channel expression contributing
to modulation of airway responsiveness in rats \cite{PMID21537608}.
Expression changes were also observed in pulmonary arterial smooth
muscle \cite{Wang_Yang_Xu_Sildenafil_inhibits_hypoxia_induced_2013}
that may relate to its antiproliferative effect
\cite{PMID18818482}
.


Happy has been the subject of several recent works
but limited progress had been made in curing her persistent
cough which is believed to be due to a 2018 heartworm infection.
Several dogs here have been treated for heartworm, including
one pregnant pit bull. While all have survived, one dog, Andy, adopted
out while still recovering in 2020. Based on second hand information, 
he was diagnosed with echo to have pulmonay arterial hypertension
and prescribed Sildenafil. Upon further investigation, that appeared
to be a good match for heartworm treated dogs with residual
damage as it may reduce right side load selectively allowing
for reverse remodelling and ultimate resolution. However, it was
not until this acute episode that it could be evaluated on Happy 
over 5 years later. 

When presented with heart failure, the prevailing standard of
care appears to include furosemide, pimobendan, and an ACE inhibitor.
Furosemide is a brute force attempt to remove excess fluid
accumulation while the ACE inhibitor is largely to remove
peripheral load. Pimobendan is largely a positive inotrope
based on calcium sensitization but also had a PDE3 inhbiting
effect possibly dilating pulmonary circulation.  In the case
of heartworm damage, which is predominantly right sided,
its not clear how beneficial peripheral dilation would be.
If there is heavy scarring in the pulmonary circulation
dilation may not be significant and the positive inotropic
effect on the right side may stimulate more maladaptive responses.
Pimbdendan is strongly contraindicated in the presence of outflow
obstructions but this warning may invoke largely concerns about
the left side. 
Pimobendan 

At least in rats with isoproterenol-induced heart damage, 
pimobendan had opposite deleterious effects on cardiac
collagen while sildenafil appeared to benefit presumably
due to difference in PDE3 and PDE5 effects
\cite{Nakata_Suzuki_Uemura_Contrasting_Effects_Inhibition_2019}.


Sildenafil is largely noted for PDE5 inhbitions and dilation
of arteries controlled by this signalling. It also has a
negative inotropic effect which may help allow reverse remodelling
by eliminating the constant mechanical loading. 





\mjmtol{
}




\newcommand{\eventtable}{
\begin{table}[H] \centering
\begin{tabular}{r|r|r|l}
\hline
Day & Serial & Date & Event \\
%\multicolumn{4}{c}{Title}\\
20604 & 0 & 2026-05-31 &last day she ate ok \\
20607 &  3 & 2026-06-03 &6 mg new furosmide started \\ 
20610 & 6 & 2026-06-06 &began eating again \\ 
20614 & 10 & 2026-06-10 &first sildenafil  5 mg  \\
20618 & 14 & 2026-06-14 &final  furosemide \\
20626 & 22 & 2026-06-22 &slowly increase  sildenafil to  6,10,10 mg  \\
20641 & 37 & 2026-07-07 &last snack until 18th, nomrinal continue  sildenafil \\
20652 & 48 & 2026-07-18 &restart snacks and regular early schedule \\ 
20652 & 48 & 2026-07-18 &not coughing much  \\ 
 & &  2026-07-23 &thought to be coughing more  \\ 
 &  & 2026-07-24 &clearly coughing a lot more  but very active suspect trachea  \\ 
\hline
\end{tabular}
\caption{}
%\label{}
\end{table}
} % eventtable

\section{Event Description  }


\eventtable

Happy began to experience symptoms consistent with
heart failure and was given expired furosemide, vetmeden, and
an ACE inhbitors which had been presdribed for her years
ago but not used. When she got to the vet, she was prescribed
prednisone, furosemide, and "Amoxi-clav." Antibiotics and furosemide
were given as prescribed but not the presnisone. She appareared
to improve to almost her prior condition although remaining a bit
slower than before and not quite as eager to eat. 
Upon request
to Happy's vet, sildenafil was prescribed at 5mg 3x per day.
She was found to do better with 6,10, and 10 mg and the furosemide
was discontinued. 
Coughing seemed to improve although since she was a bit "slow"
initially while recovering she is pretty much back to normal 
and coughing less on initial wake up and morning as well afternoon.
Does not appear to cough as much during some stresses. 
She continues to be wary of some spoon feeding suspicious for dental
pain or an unknown incident. 
As of 2026-07-08 or so she appears to continue to reduce coughing
frequency and is generally better than she has been since 2019. 
Eventually her coughing increased greatly in 2day and the sildenafil
was discontinued over concerns about collapsed traches mentioned
above. She continues to "bounce around" indicating good energy
and well being but the cough is loud and honky. 


\section{ Discussion }

Pimodendan is strongly contraindicated in the pressnce of an outflow
obstruction although some work with cats suggests it may be
safe with some obstructions
\cite{Ward_Kussin_Tropf_Retrospective_evaluation_2020}
.  This is probably true for right sided obstructions
as well as left. Scarred or hypertrophic pulmonary vasculature 
may continue to act as a right sided outflow obstruction even if
some moderate dilation can be induced. Right sided positive inotropes
would likely add stimuli for increases patholggical remodelling of both 
the heart and pulmonary system.  While pimodendam apparently does have
some PDE3 inhibiting ability, positive inotropic effects may be couterproductive. For the presumed heartworm related damage, without echo proven
PAH, sildenafil may be a better guess to allow reverse
remodelling and final resolution of heartworm maladaptations.
Vascular remodelling is an incompletely understood process that
seems to respond to details of the blood pressure waveform
\cite{Brown_Diederich_Good_Vascular_Smooth_Muscle_Remodeling_2018}
and could be expected to respond to changes in ventricular
waveform.



\begin{figure}[htb] 
\centering
\includegraphics[width=3in]{keep/Sildenafil_structure.pdf}
\caption{ 
 Sildenafil, By Mrgreen71 - Own work, CC BY-SA 3.0, {\url{https://commons.wikimedia.org/w/index.php?curid=4994287}}
Believed to be accurate but not verified.
 }
\end{figure}


\begin{figure}[htb] 
\centering
\includegraphics[width=3in]{keep/Pimobendan.pdf}
\caption{ 
Pimobendan, By Fvasconcellos - Own work, Public Domain, {\url{https://commons.wikimedia.org/w/index.php?curid=4511641}} 
Believed to be accurate but not verified.
}
\end{figure}


Happy is perhaps notable for a diet that is higher in copper and vitamin
K than most dogs.
In some cases, copper deficient diets however may reduce angioproliferation
and PAH in response to some stimuli in rats but also  stimulate 
right ventricle fibrosis  
\cite{Bogaard_Mizuno_Guignabert_Copper_dependence_angioproliferation_2011}
suggesting important tradeoffs and the importance of copper in heart health.
It is possible in some diets that  copper could be limiting in 
angiogenesis and healing from any worm related damage.
 It is thought that both of these supplementat contribute to her resilence
although the initial event and clinical course without sildenafil and
antibiotics remain unknown both of these could moderate
disease severity through several mechanisms. In the past copper appeared
to reduce coughing but as a sole focus has not entirely eliminated
Happy's cough. Essential fatty acids may have been marginal too.
Upon further investigation, vitamin D supplementation was not started
until 2020-09-19 which was well after Happy's cough free week in early 2019. Andy had received several doses.
Combined with high amino acid diet, there is some concern for excess muscl
grwoth.
Some role for vitamin D in muscle growth generally is accepted  
\cite{Polly_Tan_role_vitamin_2014} but details remain unclear and
that is probably more true in dogs than humans. 
Interestingly, salmon supplemented with high dose vitamin D
demonstrated increased cardiac experssion of a large ribosomal subunit
\cite{Gorman_Egan_AlarconLopez_Vitamin_modulates_gene_2025}
suggestive of increased ribosomal protein production.
Indeed, overexpression of vitamin D receptor was found to produce
skeletal muscle hypertrophy and ribosomal biogenesis
\cite{BASS2020101059}.
It has been characterized as promoting VSMC growth and calcification
\cite{Rebsamen_Sun_Norman_1upalpha_Dihydroxyvitamin_2002}.

Previously, Happy seemed to be sensitive to tryptophan as well
as other aromatic amino acids. 
Metabolites such as kynurenine may be elevated in PH 
\cite{PMC11068393}
but
also vasocilatory
\cite{Nagy_Nagaraj_Meinitzer_intrinsic_tryptophan_metabolite_2017}.
Tryptophan hydroxylase inhibitors are being investigrated to
block serotonin synthesis with the potential for many beneficial
effects on PAH and related remodelling
\cite{PMC11334415}
\cite{Morecroft_Dempsie_Bader_Effect_Tryptophan_Hydroxylase_2007}.

Other dietary components that may be relevant include argnine.
From the above, nitric oxide production is a good suspect to mediate 
the clinical response and arginine may be a factor.
Whlie peak amounts were as high as 250mg per sitting, which was
higher than currently, average amounts were similar or lower.
Its possible that slight elevations combined with other
unidentified factors could act with or in place of sildenafil.
One small study suggested that PAH patients metabolize arginine
differently than healthy controls
\cite{Kao_Wedes_Hsu_Arginine_metabolic_endotypes_2015}
which may be reducing NO synthesis and causing PAH although it could
also be regulated by some other cause. In Happy's case the presumed
problem is unhealed heartworm damage but mofidication of arginine metabolism
throught diet may still be possible to encourage dilation. 
In the time between 2019 and 2026, she did have up to almost
double the arginine amounts but with no remarkable changes. 


Pimobendan appears to be a very good drug in specific applications
but this has also made it trendy. While not intended for HCM
it may still be beneficial in some respects. It appeared to have
acute beneficial effects in HCM cats
\cite{Oldach_Ueda_Ontiveros_Acute_pharmacodynamic_effects_2021}
although any reactions to chronic dosing would not be observed.
It seems to be common for medicines, even blockbusters, to fix
one immediate problem while possibly creating others at least in some
populations. 
Pimobendan's most well known function is inotropy through 
calcium sensitization. It may be worth considering however if there
is some reason for the heart to reduce calcium sensitivity before
attmpting to enhance it.  Copper deficiency is one common reason
for hypertrophy and ultimate failure and there is some indication
that deficient myocytes may be prone to death with normal calcium levels
in ex vivo preparations although various channel blockers do not appear
to improve the response 
\cite{PROHASKA19991842}. 
Mice with heart specific loss of Ctr1 developed dilated cardiomyopathy
and alterations in SERCA2 related regulators 
\cite{Kim_Turski_Nose_Cardiac_Copper_Deficiency_Activates_2010}
further suggesting a link between Ca and Cu.
Prior work discusses some issues with Cu and the heart
\cite{marchywka-MJM-2024-010-0.01rg}
although not addressing Ca or pimobendan per se.
Happy's diet was notable for copper supplementation
averaging several mg per day because earlier she seemed to cough less
and act better with more copper.  Its possible that calcium desensitiation
could be adaptive to some other condition that is possible fix.  
It would be interesting to see how many pimobendan responsive dogs
do well with added copper in the right dietary context.  


\section{Conclusions}
For a dog with a history of heartworm and incompletely diagnosed
heart problems, Sildenafil may be a better guess. Pimobendan may 
even be counter productive. The nature of heartworm damage is 
not completely elucidated but perhaps reduced  signals
for patholgoical remodelling can allow for reverse remodelling
and final resolution. Its quite common in medicine to see one
thing produce some improvements in similar conditions but
lead to unappreciated damage in specific common situations.
Pimobendan appears to be trendy but may not be a good match
to the reality of many dogs' conditions. Echo and related tests
can be expensive and time consuming and not even conclusive.
"Talk to your vet about Viagra."

 

\section{Supplemental Information}

\subsection{Computer Code}


\begin{lstlisting}


\end{lstlisting}
\section{Bibliography}


\bibliography{\mjmbasename,\mjmaddbio}
\bibliographystyle{plainurl}


%%%%%%%%%%%%%%%%%%%%%%%%%%%%%%%%%%%%%%%%%%%%%%%%%%%%%%%%%%%%%%%%%%%%%%%%%%%%%
\begin{acknowledgments} 

% \input{generalack.tex}
\begin{enumerate}
\item Pubmed eutils facilities and the basic research it provides. 
\item Free software including Linux, R, LaTex  etc.
\item Thanks everyone who contributed incidental support. 
\end{enumerate}

\end{acknowledgments}

%%%%%%%%%%%%%%%%%%%%%%%%%%%%%%%%%%%%%%%%%%%%%%%%%%%%%%%%%%%%%%%%%%%%%%%%%%%%%
\clearpage
\appendix

\begin{mdpicomment}

\section{ Statement of Conflicts }
 No specific funding was used in this effort and there are no relationships
with others that could create a conflict of interest. I would like to develop
these ideas further and have obvious bias towards making them appear 
successful. Barbara Cade, the dog owner, has worked in the pet food industry
but this does not likely create a conflict. We have no interest in the
makers of any of the products named in this work.  

\end{mdpicomment}

\begin{mdpicomment}
\section{About the Authors and Facility}
This work was performed at a dog rescue run by Barbara Cade and
housed in rural Georgia.  The author of this report 
,Mike Marchywka,
has a background in electrical engineering and 
has done extensive research using free online literature sources.  
I hope to find additional people interested in critically 
examining the results and verify that they can be reproduced
effectively to treat other dogs.

\begin{comment}
\begin{figure}[htb] 
\centering
\mjmed{ picture commented out to save space in drafts...  } 
%\includegraphics[width=\picwidth]{me_on_brick.jpg}
\caption{ 
 }
\end{figure}

\end{comment}


\section{Symbols, Abbreviations and Colloquialisms}

\begin{comment}
% grep "[A-Z][A-Z]" paradox.tex | sed -e 's/[^A-Z]/\n/g' | grep "[A-Z]" | sort | uniq -c
% cat  paradox.tex | sed -e 's/  */\n/g' | grep "[A-Z][A-Z]"  | grep -v "[^A-Z]" | sort | uniq  |awk '{print $0" &   \\\\"; }'
\end{comment}


%\abbreviations{The following abbreviations are used in this manuscript:\\
%\begin{table}
\noindent
\begin{tabular}{@{}ll}
%SMVT & Sodium dependent Multi-Vitamin Transporter\\
TERM & definition and meaning   \\
\hline
%TLA & Three letter acronym\\
%LD & linear dichroism
\end{tabular} % }
%\end{table}

% https://tex.stackexchange.com/questions/5957/bibtex-entry-for-white-papers-and-technical-reports

\section{General caveats and disclaimer }
\label{appendix:caveats}

%\input{disclaimer-informal.tex}

This document was created in the hope it will be interesting to
someone including me by providing information 
about some topic that may include personal experience or a literature
review or description of a speculative theory or idea.
There is no assurance that the content of this work will be
useful for any paricular purpose. 
%In no case am I claiming to provide useful advice on any matter
%but attempting to describe events in terms of literature known
%to me. 


All statements in this document were true to the best of my knowledge
at the time they were made and every attempt is made to assure
they are not misleading or confusing. However, information provided by
others and observations that can be manipulated by unknown causes  
( "gaslighting" ) may be misleading. Any use of this information should
be preceded by validation including replication where feasible.
Errors may enter into the final work at every step from conception
and research to final editing. 
%No assurance can exist that obvious conclusions will be useful
%and may be misleading. 



Documents labelled "NOTES" or "not public" contain
substantial informal or speculative content that
may be terse and poorly edited or even sarcastic or profane.
Documents labelled as "public" have generally been edited
to be more coherent but probably have not been reviewed
or proof read. 

Generally non-public documents are labelled as such to avoid
confusion and embarassment and should be read with that understanding.


\section{Citing this as a tech report or white paper }
\label{appendix:citing}

Note: This is mostly manually entered and not assured to be error free.

This is tech report \mjmtrno. 

\begin{table}[H] \centering
\begin{tabular}{r|r|c|r}
Version & Date & Comments  &  \\
0.01 & \mjmmakedate  &  Create from empty.tex template  &  \\
-  & \today & version  \mjmversion { }   \mjmtrno  &  \\
1.0 & 20xx-xx-xx & First revision for distribution &  \\
\end{tabular}
\end{table}


Released versions,

build script needs to include empty releases.tex
\begin{table}[H] \centering
\begin{tabular}{|r|r|l|}
Version & Date & URL    \\
\hline
\expandableinput{releases.tex}
%&  &  \\
% version & date & url  \\
%.1 table & 2021-08-17& {\url{https://www.linkedin.com/posts/marchywka_draft-compare-72020-theory-with-interim-activity-6833343119203860480--wJv}} \\
%.1 table & 2021-08-17& {\url{https://www.researchgate.net/publication/353946686_Draft_table_comparing_expectations_to_recent_results_with_covid-19}} \\
%.1 table & 2021-08-17 & {\url{https://www.academia.edu/s/34e160cae9}} \\

\hline
\end{tabular}
\end{table}





% 2020-11-30 keep on same page 
%\input{bibtex2.txt}

\begin{minipage}{\linewidth}
%\input{bibtex2.txt}
%\input{bibtex3.txt}
\mjmshowbib
\end{minipage}




\begin{comment}

\end{comment}
\vspace{1cm}
Supporting files. Note that some dates,sizes, and md5's will change as this is
rebuilt.

This really needs to include the data analysis code 
but right now it is auto generated picking up things from prior
build in many cases 
\lstinputlisting{\mjmbasename.bundle_checksums}
\end{mdpicomment}
\end{document}
